\input{../../v3_rp/preamble}
\usepackage{longtable}
\usepackage{array}
\graphicspath{{../../}}
\makeatletter\def\input@path{{./}{../../}{../../v3_rp/}}\makeatother
\title{\bfseries Where Does Being Early Pay?\\
\large A catalogue of testable hypotheses for Paper B}
\author{Eric Silver}
\date{Working document, September 2026}

\begin{document}
\maketitle

\noindent Each hypothesis below says where the penalty for leading
language should be \emph{smaller} (or reversed into an advantage) and
where it should be \emph{larger}. The outcome throughout is the one the
paper already uses --- failure at the five-year proof of continued use,
plus the capital-markets and acquisition rungs where noted --- and every
test is the same comparison: the most leading fifth against the most
lagging fifth, within Nice class and year, estimated separately inside
the groups the hypothesis names and compared across them. ``Ready''
means the data exist in the repository; ``build'' means a new variable
must be constructed from data we hold; ``new data'' means a source we
do not yet have.

\section*{Where each literature looked, and who answers whom}

\begin{longtable}{>{\raggedright\arraybackslash}p{3.3cm}>{\raggedright\arraybackslash}p{5.2cm}>{\raggedright\arraybackslash}p{5.0cm}}
\toprule
Work & Setting & Cites, among these works \\
\midrule
\endhead
\citet{RobinsonFornell1985} & Surveyed consumer-goods businesses
(market share by order of entry) & no reference list deposited \\
\citet{Teece1986} & Named innovations and their imitators (theory with
cases) & no reference list deposited \\
\citet{LiebermanMontgomery1988} & Catalogue of mechanisms across
industries & Robinson \& Fornell; Teece \\
\citet{GolderTellis1993} & Histories of fifty product categories,
including diapers, video recorders, personal computers, light beer,
safety razors & no reference list deposited \\
\citet{KlepperMiller1995}; \citet{Klepper1997} & New manufactured
products through their shakeouts & no reference list deposited \\
\citet{LiebermanMontgomery1998} & Retrospective; entry timing as a
question of resources & no reference list deposited \\
\citet{KimMauborgne2004} & Cases of firms that opened new categories
(wine, fitness, circus, air travel, hotels, coffee) & cited by neither
of the later works with deposited lists \\
\citet{SuarezLanzolla2007} & Reconciliation by pace of technology and
market change & Robinson \& Fornell; Lieberman \& Montgomery 1988 and
1998; Golder \& Tellis; Teece \\
\citet{SemadeniAnderson2010} & Consulting firms' service-mark filings
(imitation of new offerings) & Lieberman \& Montgomery 1988 and 1998;
Su\'arez \& Lanzolla \\
\bottomrule
\end{longtable}

\noindent The pioneer-advantage literature answers itself in sequence;
the blue-ocean literature runs beside it. Links come from the reference
lists publishers deposit with Crossref, so a blank means the list is
not deposited, not that the citation is absent. The settings suggest
where each claim should hold if it holds anywhere: consumer packaged
goods for the pioneer claim, the named histories and shakeout-prone
manufactured goods for the counter-claim, and the blue-ocean exemplars'
own verticals (hypotheses 17--21).

\section*{A. First-mover mechanisms \citep{LiebermanMontgomery1988}}

\begin{longtable}{p{0.4cm}p{5.3cm}p{6.2cm}p{1.5cm}}
\toprule
& Claim & Test in the record & Data \\
\midrule
\endhead
1 & \emph{Technological leadership.} Early entrants protected by patents
and learning keep their lead. & Lead penalty for owners who also patent
(PatentsView match) against those who do not; expect smaller for
patenters. & ready \\
2 & \emph{Preemption of scarce assets.} Early entrants take the best
locations and inputs. & Themes built on physical sites (stores,
restaurants, hotels, clinics; classes 35, 43, 44 location vocabulary)
against footloose services; expect smaller penalty where sites matter. &
build \\
3 & \emph{Buyer switching costs.} Customers locked in by contracts stay
with the pioneer. & Descriptions naming subscriptions, memberships,
accounts or contracts against one-off goods; expect smaller penalty. &
build \\
4 & \emph{Network effects.} An installed base compounds. & Platform and
marketplace vocabulary (platform, marketplace, social network,
exchange, community); expect smaller penalty, possibly a reward. &
build \\
5 & \emph{Free-riding.} Followers copy cheaply what pioneers proved. &
Themes with low capital intensity and easy imitation (apparel, simple
consumer goods, personal services) against capital-heavy manufactured
goods; expect larger penalty. & build \\
6 & \emph{Resolution of uncertainty.} Pioneers bet before the market
settles. & Themes whose later share was volatile (surged then receded)
against themes that grew steadily; expect larger penalty in volatile
themes. & ready \\
7 & \emph{Incumbent inertia.} Established firms are slow to follow, so
the pioneer's lead lasts. & Lead penalty for debut owners against
owners with large existing portfolios filing the same leading language;
expect the penalty to fall on new firms. & ready \\
\bottomrule
\end{longtable}

\section*{B. Pace of change \citep{SuarezLanzolla2007}}

\begin{longtable}{p{0.4cm}p{5.3cm}p{6.2cm}p{1.5cm}}
\toprule
& Claim & Test in the record & Data \\
\midrule
\endhead
8 & \emph{Calm waters} (slow technology, slow market): a first-mover
advantage is durable. & Classes and eras with low turnover in the
theme mix and flat filing volume; expect the smallest penalty or a
reward. & build \\
9 & \emph{Technology leads} (fast technology, slow market): early
entrants are overtaken by later technology. & Fast theme turnover with
flat volume; expect the largest penalty. & build \\
10 & \emph{Market leads} (slow technology, fast market growth): early
entrants capture the growing demand. & Surging filing volume with a
stable theme mix (consumer categories, beverages); expect a smaller
penalty. & build \\
11 & \emph{Rough waters} (both fast): advantage only for resource-rich
entrants. & Both fast (technology classes 1995--2000); expect a large
penalty, smaller for large or funded owners. & build \\
\bottomrule
\end{longtable}

\noindent Pace of technology: year-to-year divergence in a class's theme
mix. Pace of market: growth in the class's (or theme's) filing volume.
The two cut the class-years into the four quadrants.

\section*{C. Appropriability \citep{Teece1986}}

\begin{longtable}{p{0.4cm}p{5.3cm}p{6.2cm}p{1.5cm}}
\toprule
& Claim & Test in the record & Data \\
\midrule
\endhead
12 & \emph{Tight appropriability}: where innovations can be protected,
the innovator profits. & Pharmaceuticals and medical devices (classes 5,
10) and patenting owners; expect smaller penalty. & ready \\
13 & \emph{Weak appropriability, generic complementary assets}:
imitators win. & Services and software with no patenting; expect larger
penalty. & ready \\
14 & \emph{Specialized complementary assets held by incumbents}:
incumbents who own distribution capture the innovation. & In
distribution-heavy consumer goods (food, beverages, cosmetics), the
interaction of lead with owner size: leading small entrants fail,
leading large incumbents survive. & ready \\
\bottomrule
\end{longtable}

\section*{D. Pioneer disadvantage \citep{GolderTellis1993} and
shakeouts \citep{KlepperMiller1995, Klepper1997}}

\begin{longtable}{p{0.4cm}p{5.3cm}p{6.2cm}p{1.5cm}}
\toprule
& Claim & Test in the record & Data \\
\midrule
\endhead
15 & \emph{Early leaders beat pioneers}: the firms remembered as
pioneers usually entered a few years after the first entrants. & Within
each theme's arrival in a class, the first year's filers against those
two to five years in; expect the later cohort to survive better. &
ready \\
16 & \emph{Survivor bias}: pioneer advantage appears only among
survivors. & The lead gradient among all applications, among
registrations, and among marks that passed the proof; expect the
gradient to fade as the sample is conditioned on survival. & ready \\
17 & \emph{Pioneer advantage in consumer packaged goods} (the finding
Golder and Tellis dispute). & Consumer packaged goods classes (3, 29,
30, 32, 33) against industrial goods; the pioneer literature predicts
opposite signs. & ready \\
18 & \emph{Early entrants win shakeouts} in manufactured products with
scale economies. & Manufactured goods classes (7, 9, 12) against
services; and, within surging themes, early joiners against late
joiners (already: no order effect). & ready \\
19 & \emph{The cited cases}: disposable diapers, video recorders,
personal computers, light beer, safety razors. & Lead penalty within
the verticals those cases come from (paper and sanitary goods, consumer
electronics, beer, razors and blades). & ready \\
\bottomrule
\end{longtable}

\section*{E. Blue ocean \citep{KimMauborgne2004}}

\begin{longtable}{p{0.4cm}p{5.3cm}p{6.2cm}p{1.5cm}}
\toprule
& Claim & Test in the record & Data \\
\midrule
\endhead
20 & \emph{Reconstructing market boundaries}: value comes from combining
attributes from different industries. & Filings carrying a theme
established in another class into their own (new-to-class), and
filings whose leading language is a new pairing of existing themes;
expect a reward. & ready \\
21 & \emph{The exemplars' verticals}: wine ([yellow tail]), fitness
(Curves), live entertainment (Cirque du Soleil), air travel (Southwest,
NetJets), hotels, coffee, home improvement. & Lead penalty within each
vertical's class and themes; and where the exemplar's own filings fall
in the lead distribution when they are in the scored years. & ready \\
22 & \emph{Creating new demand} (non-customers) rather than taking share.
& Themes whose growth came from first-time filers rather than
established owners; expect a smaller penalty for the early entrants of
such themes. & build \\
\bottomrule
\end{longtable}

\section*{F. Kinds of new}

\begin{longtable}{p{0.4cm}p{5.3cm}p{6.2cm}p{1.5cm}}
\toprule
& Claim & Test in the record & Data \\
\midrule
\endhead
23 & \emph{New inventions} reward pioneers who define the category. &
Filings carrying a theme new to the whole corpus (too few at 500 themes;
use curated invention vocabularies: 3D printing, drones, gene therapy,
electric vehicles). & build \\
24 & \emph{Fads} punish late entry, and possibly early entry too. &
Themes that surged and then fell back below their pre-surge share
within five years, against surges that held. & ready \\
25 & \emph{General-purpose technologies} (internet, mobile, cloud, AI)
spread across many industries; early adopters in each industry may
win. & Vocabularies appearing in ten or more classes; the lead penalty
inside each adopting class, by years since the vocabulary reached it
(the internet result generalized). & ready \\
26 & \emph{New to the class versus new to the world.} Importing an
established theme is safer than originating one. & New-to-class against
new-to-corpus themes, and both against established themes. & ready \\
27 & \emph{Groups of emergent consumer themes} (plant-based milk,
kombucha, cold brew, hard seltzer, CBD, energy drinks). & Pooled into
one ``emergent consumer category'' group, with each category's early
entrants compared with its later entrants and with the class. & ready \\
\bottomrule
\end{longtable}

\section*{G. Who and where}

\begin{longtable}{p{0.4cm}p{5.3cm}p{6.2cm}p{1.5cm}}
\toprule
& Claim & Test in the record & Data \\
\midrule
\endhead
28 & \emph{Clusters speed imitation.} Where a theme's early filers are
geographically concentrated (a Silicon Valley or Silicon Alley),
followers learn fastest. & Co-location index of the theme's filers
in a five-year window, net of its class (see the clustering measure
below); expect a larger penalty in concentrated themes. & ready (ZIP) \\
29 & \emph{Distance from the cluster.} Early filers far from a theme's
hubs are shielded from local imitation, or cut off from their
resources. & Within themes that have a hub, filers in a hub against
filers outside; the sign is the question. & ready (ZIP) \\
30 & \emph{Foreign pioneers.} Foreign firms bringing a theme into the US
market face local imitators. & Foreign-domiciled and Madrid filers
against US filers. & ready \\
31 & \emph{Resources absorb the penalty.} & Counsel-represented and
large-portfolio owners against self-filed debut owners. & ready \\
32 & \emph{Speculation.} Intent-to-use filings claim a market before
entering it. & Intent-to-use against use-based filings; expect a larger
penalty for speculative leading filings. & ready \\
33 & \emph{Regulation as a barrier.} Approval requirements protect the
first entrant. & Regulated classes (pharmaceuticals, medical devices,
alcohol, tobacco, financial services) against unregulated ones. &
ready \\
34 & \emph{Timing of entry in the cycle.} Early entry in a boom is
punished; in a bust it is rewarded. & The era swing, extended by the
reversal decomposition to every class and theme. & ready \\
\bottomrule
\end{longtable}

\section*{Notes on design}

Many of these groupings overlap, and several are proxies (a vocabulary
for ``switching costs'' is not a measured switching cost). The battery
should be run with one pre-stated prediction per hypothesis, a common
estimator, owner-clustered errors, and a correction for testing thirty
comparisons at once; the paper would then report where the evidence
lines up with each literature, not a search for significance. Owner
locations are the applicant's ZIP code as filed, re-parsed from the
annual case-file archives and placed at Census ZIP centroids (99.8\% of
US owners placed).

\section*{What the battery shows}

Across all registrations of 2002--2018 the most leading fifth fails the
five-year proof 1.39 points more often than the most lagging fifth. Six
of the thirty pre-stated contrasts survive adjustment for multiple
testing; four of the six run against the prediction.

\emph{The penalty is concentrated where a class's filing volume is
growing.} Where both the class's theme mix and its filing volume changed
fast, the penalty is 3.2 points; where only volume grew, 1.9; in calm
class-years, 0.6; where only the mix turned over, $-0.2$. The pace-of-change
account predicts the most durable advantage where the market grows and the
technology is stable \citep{SuarezLanzolla2007}; the record puts the
penalty there instead. This is what the crowding reading of lead
predicts: when entrants pour into a class, the language they pour into is
the language that fails.

\emph{Platform and network offerings carry the largest penalty.} Filings
describing marketplaces, platforms, social networks or peer-to-peer
services have a penalty of 7.2 points, against 1.2 for all other filings.
Network effects are listed as a source of first-mover advantage
\citep{LiebermanMontgomery1988}; in the record, early entrants to
networked categories fail more often than late ones.

\emph{In consumer packaged goods the penalty falls on established firms,
not new ones.} Among owners with 25 or more earlier filings, leading
filings fail 4.9 points more often; among first-time filers, 0.9 points
less often. The complementary-assets account predicts the reverse
\citep{Teece1986}: incumbents with distribution should hold what small
pioneers open.

\emph{Manufactured goods carry a larger penalty than services} (2.3
against 1.1 points), the reverse of the shakeout account, in which early
manufacturers are the likeliest survivors \citep{KlepperMiller1995,
Klepper1997}.

\emph{Counsel absorbs most of the penalty.} Self-filed leading filings
fail 4.2 points more often than self-filed lagging ones; counsel-represented,
0.6.

\emph{Two contrasts fall just short of the adjustment} (Holm $p = 0.065$):
the penalty is larger in themes whose share swings from year to year (2.1
against 0.9), as the uncertainty account predicts, and larger for
site-based offerings such as restaurants, hotels and clinics (2.7 against
1.4), against the preemption account.

\emph{Pioneers of a new vocabulary fail, and so do its early followers.}
Filings using a new technology's or product's vocabulary in its first two
years in a class fail 12.8 points more often than other filings of the
same class and year; filings two to five years in, 11.8 points more
often; filings six or more years in, 6.0. The first two do not differ
($t = 0.5$). Golder and Tellis's pioneers failed; here the early followers
failed at the same rate.

\emph{Within verticals the sign differs.} Leading filings in beer fail
6.2 points less often than lagging ones, and in live entertainment 5.6
points less often; in coffee they fail 3.2 points more often and in
fitness 2.9. These vertical estimates are not adjusted for multiple
testing.

\emph{The blue-ocean exemplars' own marks do not lead their classes.}
The mean lead fifth of the registrations of [yellow tail]'s owner,
Cirque du Soleil, Southwest Airlines, NetJets, Accor, Starbucks and Home
Depot lies between 2.7 and 3.2, the middle of the distribution; only
Curves's marks lean leading (3.8). These are the firms' marks of
2002--2018, mostly filed years after each created its category.

\emph{Geography changes where filers fail, not the lead penalty.} A theme's
co-location with itself does not change its penalty (1.8 in the most
concentrated third, 2.3 in the most dispersed). Within themes that have a
hub, filers in the hub fail 1.2 points more often than filers outside it
(s.e.\ 0.45), whatever their lead, and their lead penalty is the same.

\emph{Nothing distinguishes} patenting owners, pharmaceuticals and
devices, consumer against industrial goods, themes imported from another
class, themes grown by first-time filers, invention vocabularies, fads
against surges that held, vocabulary new to its class against new to the
corpus, foreign owners, intent-to-use filings, regulated classes, or
boom against bust years. Several of these groups are small (invention
vocabularies 12,612 registrations; general-purpose-technology filings in
their first three years in a class, 1,166).


\section*{Results}
Penalty = failure rate at the five-year proof of the most leading fifth minus the most lagging fifth, in percentage points (all registrations 2002--2018: +1.39, s.e.\ 0.17). Each row's two penalties come from one model with the lead fifths interacted with the group; the difference is the first minus the second. Predicted: the sign of that difference stated in the catalogue (? = no prediction). Holm $p$ adjusts across all pre-stated contrasts. Agrees: whether the sign matches the prediction, given only where Holm $p < 0.10$.

{\small
\begin{longtable}{>{\raggedright\arraybackslash}p{0.7cm}>{\raggedright\arraybackslash}p{4.6cm}>{\raggedright\arraybackslash}p{2.9cm}rrrcr}
\toprule & Groups compared & Penalty (pp) & Diff. & $t$ & Holm $p$ & Pred. & Agrees \\ \midrule \endhead
1 & Owners who patent vs those who do not & patenter +1.10; no patents +1.50 & -0.39 & -0.6 & 1.000 & - & -- \\
2 & Site-based offerings vs others & site-based +2.68; other +1.40 & +1.28 & 3.0 & 0.065 & - & no \\
3 & Contract / subscription offerings vs others & switching costs -0.32; other +1.51 & -1.83 & -2.7 & 0.159 & - & -- \\
4 & Platform / network offerings vs others & network +7.17; other +1.17 & +6.00 & 5.8 & 0.000 & - & no \\
5 & Easy-to-imitate goods vs capital-heavy goods & easy to imitate -0.01; capital heavy +1.53 & -1.54 & -2.5 & 0.232 & + & -- \\
6 & Volatile themes vs steady themes (top vs bottom third of share volatility) & volatile +2.10; steady +0.93 & +1.18 & 3.0 & 0.065 & + & yes \\
7 & Debut owners vs established owners (25+ prior filings) & debut +1.73; established +1.36 & +0.37 & 0.6 & 1.000 & + & -- \\
8 & Pace: technology leads vs calm & technology leads -0.24; calm +0.55 & -0.78 & -2.2 & 0.595 & + & -- \\
10 & Pace: market leads vs technology leads & market leads +1.92; technology leads -0.24 & +2.15 & 5.4 & 0.000 & - & no \\
11 & Pace: rough vs calm & rough +3.21; calm +0.55 & +2.66 & 7.0 & 0.000 & + & yes \\
12 & Pharmaceuticals and medical devices (classes 5, 10) vs all other classes & pharma/devices +2.88; other +1.31 & +1.57 & 1.3 & 1.000 & - & -- \\
13 & Services (classes 35-45) vs goods & services +1.12; goods +1.56 & -0.43 & -1.3 & 1.000 & + & -- \\
14 & Consumer packaged goods: debut owners vs established owners & debut -0.86; established +4.89 & -5.75 & -4.6 & 0.000 & + & no \\
15 & Pioneers of a new vocabulary in a class vs early followers (excess failure over same class-year) & pioneer --; early follower -- & +0.95 & 0.5 & 1.000 & + & -- \\
17 & Consumer packaged goods vs industrial goods & CPG +1.57; industrial +0.96 & +0.61 & 0.9 & 1.000 & - & -- \\
18 & Manufactured goods (7, 9, 12) vs services (35-45) & manufactured +2.33; services +1.08 & +1.25 & 3.2 & 0.034 & - & no \\
20 & Theme imported from another class vs native theme & imported +1.92; native +1.35 & +0.56 & 1.3 & 1.000 & - & -- \\
22 & Themes growing through first-time filers vs through established filers (top vs bottom third) & new demand +1.05; established demand +0.33 & +0.72 & 1.6 & 1.000 & - & -- \\
23 & Invention vocabularies vs all other filings & invention +3.33; other +1.38 & +1.95 & 1.1 & 1.000 & - & -- \\
24 & Fads vs surges that held & fad (surge given back) +1.41; surge held +2.87 & -1.46 & -0.5 & 1.000 & + & -- \\
25 & General-purpose technology filings: first 3 years in a class vs 6+ years & early adopter -2.30; late adopter +1.86 & -4.16 & -0.8 & 1.000 & - & -- \\
26 & Vocabulary new to its class vs new to the corpus (excess failure, first five years) & new to class +11.09; new to corpus +9.46 & +1.63 & 0.6 & 1.000 & - & -- \\
27 & Emergent consumer categories: first three years vs later (excess failure) & first 3 years +12.62; later +12.63 & -0.01 & -0.0 & 1.000 & ? &  \\
28 & Geographically concentrated themes vs dispersed (top vs bottom third of co-location L) & concentrated +1.80; dispersed +2.31 & -0.51 & -1.1 & 1.000 & + & -- \\
29 & Within themes that have a hub: filers in a hub vs filers outside & in hub +2.81; outside hub +3.08 & -0.27 & -0.4 & 1.000 & ? &  \\
30 & Foreign-domiciled owners vs US owners & foreign (not China) +0.63; US +1.65 & -1.02 & -1.9 & 0.987 & + & -- \\
31 & Self-filed vs counsel-represented & self-filed +4.20; counsel +0.59 & +3.61 & 7.7 & 0.000 & + & yes \\
32 & Intent-to-use vs use-based filings & intent to use +1.43; in use +1.36 & +0.07 & 0.2 & 1.000 & + & -- \\
33 & Regulated classes (5, 10, 33, 34, 36) vs others & regulated +1.96; unregulated +1.31 & +0.65 & 1.0 & 1.000 & - & -- \\
34 & Filed in a boom (2004-07, 2014-17) vs a bust (2001-02, 2008-10) & boom +0.92; bust +1.01 & -0.09 & -0.3 & 1.000 & + & -- \\
\bottomrule\end{longtable}}

Rows 15, 26 and 27 report excess failure (pp) over other filings in the same class and registration year, not a lead penalty.

\paragraph{Pace of change (hypotheses 8--11).} Penalty by quadrant: calm +0.55 (n = 830,672); technology leads -0.24 (n = 718,425); market leads +1.92 (n = 707,755); rough +3.21 (n = 847,682).
\paragraph{Golder--Tellis cases (hypothesis 19).} Penalty within each: diapers +0.36 (s.e.\ 1.97, n = 9,788); video recorders +3.65 (s.e.\ 3.47, n = 10,306); personal computers +2.19 (s.e.\ 1.54, n = 27,690); beer -6.16 (s.e.\ 1.33, n = 39,290); razors +0.34 (s.e.\ 2.93, n = 4,090).
\paragraph{Blue-ocean exemplar verticals (hypothesis 21).} Penalty within each: wine +1.29 (s.e.\ 0.94, n = 49,428); fitness +2.93 (s.e.\ 1.06, n = 47,894); live entertainment -5.58 (s.e.\ 1.50, n = 26,045); air travel +1.21 (s.e.\ 2.94, n = 5,729); hotels -1.67 (s.e.\ 1.47, n = 24,879); coffee +3.16 (s.e.\ 1.07, n = 38,624); home improvement +1.10 (s.e.\ 1.96, n = 13,176).
\paragraph{The exemplar firms' own marks.} Registrations 2002--2018 by owner; mean lead fifth (1 = most lagging, 5 = most leading), share in the most leading fifth, and proof failure rate: {[}yellow tail{]} (Casella) 109 marks, fifth 3.0, 23\% leading, 48\% failed; Curves 65 marks, fifth 3.8, 42\% leading, 51\% failed; Cirque du Soleil 64 marks, fifth 2.7, 22\% leading, 34\% failed; Southwest Airlines 103 marks, fifth 2.7, 12\% leading, 42\% failed; NetJets 9 marks, fifth 2.7, 11\% leading, 100\% failed; Accor (Formule 1) 284 marks, fifth 2.9, 19\% leading, 45\% failed; Starbucks 239 marks, fifth 3.2, 30\% leading, 42\% failed; Home Depot 346 marks, fifth 2.9, 18\% leading, 17\% failed.
\paragraph{Being in a hub.} Within themes that have a hub, filers in the hub fail +1.19 points more often than filers outside it (s.e.\ 0.45), whatever their lead.
\section*{The clustering measure}
For a group of $N$ US filers placed at their ZIP centroids, the co-location index is the average over pairs of distinct filers of $k(d) = \max(0, 1 - d/200\text{ miles})$: two filers in the same place score 1, two 100 miles apart 0.5, two more than 200 miles apart 0. It is the chance that two filers drawn at random are neighbours, weighted by how near, and it allows any number of hubs. Filers crowd into populous places whatever they sell, so the index is set against the same index $B$ for all filers in the class and years: $L = (C - B)/(1 - B)$, which is 0 for a group placed like its class and 1 when every filer is in one place. This is the distance-based comparison of \citet{DurantonOverman2005} summarized at one bandwidth. A hub is a place where at least 10\% of the group's filers are nearby (kernel-weighted) and at least 1.2 times the class's share; hubs are picked greedily and named by the commonest owner city within 50 miles. For the battery the groups are each class's filers of one primary theme in a five-year window around the filing year, with at least 30 US filers.

{\small\begin{longtable}{>{\raggedright\arraybackslash}p{5.2cm}rrrr>{\raggedright\arraybackslash}p{4.6cm}}
\toprule Group & $N$ & $C$ & $B$ & $L$ & Hubs (share nearby vs class) \\ \midrule \endhead
internet vocabulary, classes 9/35/38/42, 1998-2001 & 60,903 & 0.075 & 0.062 & 0.014 & San Francisco, CA 13\% vs 10\% \\
semiconductors, class 9, 1995-2005 & 4,517 & 0.254 & 0.064 & 0.203 & Santa Clara, CA 48\% vs 12\% \\
wine, class 33, 2000-2010 & 18,799 & 0.159 & 0.108 & 0.056 & Napa, CA 35\% vs 26\% \\
country music, class 41, 1995-2010 & 163 & 0.106 & 0.067 & 0.041 & Nashville, TN 28\% vs 1\% \\
casino gaming, class 41, 1995-2010 & 4,662 & 0.156 & 0.067 & 0.095 & Las Vegas, NV 37\% vs 2\% \\
biotechnology, classes 1/5/42, 1995-2010 & 2,695 & 0.090 & 0.060 & 0.031 & Philadelphia, PA 16\% vs 13\%; Foster City, CA 14\% vs 7\%; New Haven, CT 13\% vs 11\% \\
social networking, classes 38/42/45, 2004-2012 & 14,041 & 0.088 & 0.057 & 0.033 & New York, NY 18\% vs 14\%; San Francisco, CA 15\% vs 8\%; Los Angeles, CA 13\% vs 9\% \\
online advertising, classes 35/42, 2005-2015 & 12,028 & 0.079 & 0.054 & 0.027 & New York, NY 17\% vs 14\%; San Francisco, CA 14\% vs 7\%; Los Angeles, CA 12\% vs 9\% \\
software as a service, class 42, 2008-2016 & 20,783 & 0.064 & 0.057 & 0.008 & San Francisco, CA 15\% vs 10\% \\
hip hop, classes 9/25/41, 1995-2010 & 507 & 0.122 & 0.068 & 0.058 & New York, NY 29\% vs 16\% \\
\bottomrule\end{longtable}}


\begin{thebibliography}{99}

\bibitem[Golder \& Tellis(1993)]{GolderTellis1993}
Golder, P.~N., \& Tellis, G.~J. (1993). Pioneer advantage: Marketing logic or
marketing legend? \emph{Journal of Marketing Research}, 30(2), 158--170.

\bibitem[Kim \& Mauborgne(2004)]{KimMauborgne2004}
Kim, W.~C., \& Mauborgne, R. (2004). Blue ocean strategy. \emph{Harvard
Business Review}, 82(10), 76--84.

\bibitem[Klepper(1997)]{Klepper1997}
Klepper, S. (1997). Industry life cycles. \emph{Industrial and Corporate
Change}, 6(1), 145--182.

\bibitem[Klepper \& Miller(1995)]{KlepperMiller1995}
Klepper, S., \& Miller, J.~H. (1995). Entry, exit, and shakeouts in the
United States in new manufactured products. \emph{International Journal
of Industrial Organization}, 13(4), 567--591.

\bibitem[Robinson \& Fornell(1985)]{RobinsonFornell1985}
Robinson, W.~T., \& Fornell, C. (1985). Sources of market pioneer advantages
in consumer goods industries. \emph{Journal of Marketing Research}, 22(3),
305--317.

\bibitem[Lieberman \& Montgomery(1988)]{LiebermanMontgomery1988}
Lieberman, M.~B., \& Montgomery, D.~B. (1988). First-mover advantages.
\emph{Strategic Management Journal}, 9(S1), 41--58.

\bibitem[Lieberman \& Montgomery(1998)]{LiebermanMontgomery1998}
Lieberman, M.~B., \& Montgomery, D.~B. (1998). First-mover (dis)advantages:
Retrospective and link with the resource-based view. \emph{Strategic
Management Journal}, 19(12), 1111--1125.

\bibitem[Semadeni \& Anderson(2010)]{SemadeniAnderson2010}
Semadeni, M., \& Anderson, B.~S. (2010). The follower's dilemma:
Innovation and imitation in the professional services industry.
\emph{Academy of Management Journal}, 53(5), 1175--1193.

\bibitem[Su\'arez \& Lanzolla(2007)]{SuarezLanzolla2007}
Su\'arez, F.~F., \& Lanzolla, G. (2007). The role of environmental
dynamics in building a first mover advantage theory. \emph{Academy of
Management Review}, 32(2), 377--392.

\bibitem[Teece(1986)]{Teece1986}
Teece, D.~J. (1986). Profiting from technological innovation:
Implications for integration, collaboration, licensing and public
policy. \emph{Research Policy}, 15(6), 285--305.

\end{thebibliography}

\end{document}
